% Pregunta 212 (2021-julio-OEP2020, n.º 17): IS-LM-BP con tipo de cambio fijo y BP menos
% inclinada que la LM; política monetaria expansiva. Rectas a escala:
%   LM: i = 2,5 + 1,2 (Y - 2,5)   BP: i = 2,5 + 0,5 (Y - 2,5)   IS: i = 2,5 - 0,8 (Y - 2,5)
%   LM': i = 2,5 + 1,2 (Y - 3,75)
%   E0 = (2,5; 2,5)   E1 = IS∩LM' = (3,25; 1,9), por debajo de la BP (2,875): déficit
% Autor: Víctor Gutiérrez Marcos
\documentclass[tikz,border=6pt]{standalone}
\input{tikz-tcee}
\begin{document}
\begin{tikzpicture}[tcee,x=1.15cm,y=1.15cm]
  \ejes{$Y$}{$i$}{6}{5}
  % Curvas iniciales
  \draw[curva3] (0.3,1.4) -- (5.4,3.95) node[etiqueta,right,text=tceeverde] {$BP$};
  \draw[curva2] (0.4,4.18) -- (5.3,0.26) node[etiqueta,right,text=tceeazul] {$IS$};
  \draw[curva1] (0.8,0.46) -- (4.2,4.54) node[etiqueta,above,text=tceemorado] {$LM$};
  % Lo que cambia
  \draw[curvanueva,desplazada] (2.05,0.46) -- (5.45,4.54) node[etiqueta,above,text=tceerojo] {$LM'$};
  \draw[flecha] (4.95,4.0) -- (3.68,4.0);
  \node[etiqueta,font=\footnotesize,text=tceerojo,align=right,anchor=east] at (3.38,4.0) {venta de\\reservas};
  % Proyecciones
  \proyeccion{(2.5,2.5)}{$Y_0$}{$i_0$}
  \proyeccion{(3.25,1.9)}{$Y_1$}{$i_1$}
  % Puntos
  \node[punto] at (2.5,2.5) {};
  \node[etiqueta,above left] at (2.5,2.5) {$E_0$};
  \node[punto] at (3.25,1.9) {};
  \node[etiqueta,anchor=west] at (3.36,1.95) {$E_1$};
  \draw[flecha] (3.12,2.0) to[bend left=25] (2.62,2.4);
  \node[etiqueta,font=\footnotesize,text=tceerojo,align=left,anchor=west] at (4.0,2.05) {déficit\\(bajo la $BP$)};
\end{tikzpicture}
\end{document}
